\documentclass[svgnames]{article}     % use "amsart" instead of "article" for AMSLaTeX format
%\geometry{landscape}                 % Activate for rotated page geometry

%\usepackage[parfill]{parskip}        % Activate to begin paragraphs with an empty line rather than an indent

\usepackage{graphicx}                 % Use pdf, png, jpg, or eps§ with pdflatex; use eps in DVI mode

%maths                                % TeX will automatically convert eps --> pdf in pdflatex
\usepackage{amssymb}
\usepackage{amsmath}
\usepackage{esint}
\usepackage{geometry}

% Inverting Color of PDF
\usepackage{xcolor}
\pagecolor[rgb]{0.19,0.19,0.19}
\color[rgb]{0.77,0.77,0.77}

%noindent
\setlength\parindent{0pt}

%pgfplots
\usepackage{pgfplots}

%images
%\graphicspath{{ }}                   % Activate to set a image directory

%tikz
\usepackage{pgfplots}
\pgfplotsset{compat=1.15}
\usepackage{comment}
\usetikzlibrary{arrows}
\usepackage[most]{tcolorbox}

%Figures
\usepackage{float}
\usepackage{caption}
\usepackage{lipsum}

\usepackage{listings}
\usepackage{color}

\definecolor{dkgreen}{rgb}{0,0.6,0}
\definecolor{gray}{rgb}{0.5,0.5,0.5}
\definecolor{mauve}{rgb}{0.58,0,0.82}

\lstset{frame=tb,
  language=Java,
  aboveskip=3mm,
  belowskip=3mm,
  showstringspaces=false,
  columns=flexible,
  basicstyle={\small\ttfamily},
  numbers=none,
  numberstyle=\tiny\color{gray},
  keywordstyle=\color{blue},
  commentstyle=\color{dkgreen},
  stringstyle=\color{mauve},
  breaklines=true,
  breakatwhitespace=true,
  tabsize=3
}

\lstset{language=Python}

\tcbset {
  base/.style={
    arc=0mm,
    bottomtitle=0.5mm,
    boxrule=0mm,
    colbacktitle=black!10!white,
    coltitle=black,
    fonttitle=\bfseries,
    left=2.5mm,
    leftrule=1mm,
    right=3.5mm,
    title={#1},
    toptitle=0.75mm,
  }
}

\definecolor{brandblue}{rgb}{0.34, 0.7, 1}
\newtcolorbox{mainbox}[1]{
  colframe=brandblue,
  base={#1}
}

\newtcolorbox{subbox}[1]{
  colframe=black!30!white,
  base={#1}
}

\title{Single Systems -- Lesson 01}
\author{Deval Deliwala}
%\date{}                              % Activate to display a given date or no date

\begin{document}
\maketitle
%\section{}
%\subsection{}
\tableofcontents                     % Activate to display a table of contents
\newpage

\paragraph{Notation} \mbox{} \\

 -- $X$ refers to the system being considered \\
 -- $\Sigma$ refers to the set of classical states of $X$ \\

Here are a few examples: 

\begin{itemize}
  \item[-] if $X$ is a bit, $\Sigma = {0, 1}$ -- the \textit{binary alphabet}
  \item[-] if $X$ is a six-sided die, $\Sigma = {1,2,3,4,5,6}$ 
  \item[-] if $X$ is an electric fan switch, $\Sigma = {\text{high, medium, low, off}} $
\end{itemize} 

\section{Classical Information}

\subsection{Classical States \& Probability Vectors}

In Quantum Computing (QC) our knowledge of $X$ is uncertain. We thus represent our knowledge of the classical state of $X$ by assigning \textit{probabilities} to each classical state resulting in a  \textit{probabilistic state}. \\

For example, suppose $X$ is a bit. In this case, based on our past experience or what we know about $X$, there is a $3/4$ chance its classical state is 0 and a $1/4$ chance it's 1. Therefore,

\[
\text{Pr} (X = 0) = \frac{3}{4} \qquad \text{Pr} (X = 1) = \frac{1}{4} 
\] \vspace{5px}

We can represent this more succinctly with a column vector: 

\[
\begin{pmatrix}
  \frac{3}{4}\\[5px]
  \frac{1}{4}
\end{pmatrix}
\] \vspace{5px}

The probability of the bit being 0 is placed at the top; the probability of the bit being 1 is placed at the bottom. We can represent any probabilistic state via a column vector satisfying two properties: \\ 

\begin{mainbox}{Probabilistic Vector Requirements}
  \begin{itemize}
    \item[1.] All entries of the vector are \textit{nonnegative numbers}
    \item[2.] The sum of the entries is equal to 1
  \end{itemize}
\end{mainbox}

\subsection{Measuring Probabilistic States}

Intuitively, we can never ``see" a system in a probabilistic state; a measurement always yields exactly one of the allowed states. \\

Measuring changes our knowledge of the system, and therefore changes the probabilistic state we associate with the system. If we recognize that $X$ is in the classical state $a \in \Sigma$, then the new probability vector representing our knowledge of $X$ becomes a vector having 1 in the entry corresponding to $a$ and 0 for all other entries. 

\[
\begin{pmatrix}
  0.3 \\
  0.1 \\
  3.14 \\
  2.72 \\
  \vdots 
\end{pmatrix} 
\quad \rightarrow \quad \boxed{\text{Measurement}} \quad \rightarrow \quad
\begin{pmatrix}
  0 \\
  0 \\
  1 \\
  0 \\
  \vdots
\end{pmatrix}
\] \vspace{5px}

\subsection{Standard Basis Vectors}

We can define any probabilistic state vector as a \textit{linear combination} of standard basis vectors. For example, assuming the system we have in mind is a bit, the standard basis vectors are given by 

\begin{subbox}{Computational Basis States}
  \[
  |0\rangle = \begin{pmatrix}
    1 \\ 0  
  \end{pmatrix} \quad \text{and} \quad |1 \rangle = \begin{pmatrix}
    0 \\ 1
  \end{pmatrix}.
\]
\end{subbox}

For example, we have

\[
\begin{pmatrix}
  \frac{3}{4}\\[5px] \frac{1}{4} 
\end{pmatrix} = \frac{3}{4} |0\rangle + \frac{1}{4}|1\rangle
\] \vspace{5px}
\subsection{Operations}

\paragraph{Deterministic Operations} \mbox{} \\

Deterministic Operations transform each classical state $a \in \Sigma$ into $f(a)$ for some function f of the form $f: \Sigma \rightarrow \Sigma$. \\

For example, if  $\Sigma = {0,1}$, there are four functions of the form $f_1, f_2, f_3, f_4$ which can be represented as follows: 

\[
f_1(0) = 0 \quad f_1(1) = 0 
\] 
\[ f_2(0) = 0 \quad f_2(1) = 1 \]
\[f_3(0) = 1 \quad f_3(1) = 0 \]
\[ f_4(0) = 1 \quad f_4(1) = 1 \] \vspace{5px}

The first and last of these functions are  \textit{constant}, where the output remains constant regardless of input. The middle two are \textit{balanced} where the two possible output values occur the same \# of times. The function $f_2$ is the \textbf{identity function} where $f(a) = a$ for $a \in \Sigma$. The function $f_3$ is the  $NOT$ function where each input is flipped for an output. \\

Every deterministic operation on probabilistic states can be represented as a matrix, where

\[
M |a\rangle = |f(a)\rangle 
\] \vspace{5px}
for every $a \in \Sigma$. Such a matrix always exists and is unique. For the above constant and balanced functions: 

\[
M_1 = \begin{pmatrix}
  1 & 1 \\ 0 & 0
\end{pmatrix}, \quad M_2 = \begin{pmatrix}
  1 & 0 \\ 0 & 1
\end{pmatrix}, \quad M_3 = \begin{pmatrix}
  0 & 1 \\ 1 & 0
\end{pmatrix}, \quad M_4 = \begin{pmatrix}
  0 & 0 \\ 1 & 1
\end{pmatrix} \] \vspace{5px}

Deterministic matrix operations always have exactly one 1 in each column, and 0 for all other entries. \\

Let us denote $\langle a |$ the \textit{row} vector having a 1 in the entry corresponding to $a$ and 0 for all other entries, for each $a \in \Sigma$. \\

For example if $\Sigma = {0, 1}$, 

\[
  \langle 0 | = ( 1 \quad 0 ) \quad \text{and} \quad \langle 1 | = (0 \quad 1) \] \vspace{5px}

If we perform matrix multiplication on a column vector defined by $|b\rangle$ and a row vector defined by $\langle a |$, we obtain a square matrix having a 1 in the entry corresponding to the $(b, a)$ location in the matrix and 0 everywhere else. For example, 

\[
  |0\rangle\langle 1 | = \begin{pmatrix}
    0 & 1 \\ 0 & 0
\end{pmatrix}.  \] \vspace{5px}
Using this notation, we can express $M$ corresponding to the function $f$ as 

\[
  M = \sum_{a \in \Sigma} | f(a)\rangle\langle a |
\] \vspace{5px}
If we switch the order of multiplication -- $\langle a || b\rangle$, we obtain a 1 x 1 scalar. For the sake of tidiness we write the product as $\langle a | b \rangle$. We will later define $\langle a | b \rangle$ as the \textit{inner product} of $a $ and $b$.  

 \[
\langle a | b \rangle = \begin{cases}
  1 \quad &a=b \\ 0 \quad &a\neq b
\end{cases} \] \vspace{5px}


\subsection{Probabilistic Operations \& Stochastic Matrices}

In addition to deterministic operations, we have \textit{probabilistic
operations.}

For example, consider an operation on a bit where, if the classical state of
the bit is 0, it is left alone, and if the classical state of the bit is 1, it
is flipped to 0 with probability $\frac{1}{2}$. This operation is represented
by the matrix

\[
\begin{pmatrix}
  1 & \frac{1}{2} \\[2px] 
  0 & \frac{1}{2}
\end{pmatrix}
\] \vspace{5px}

We can check this matrix does the correct thing by multiplying with the
standard basis vectors: 

\[
\begin{pmatrix}
  1 & \frac{1}{2} \\[2px]
  0 & \frac{1}{2} 
\end{pmatrix} 
\begin{pmatrix}
  1 \\[2px] 0
\end{pmatrix} = \begin{pmatrix}
1 \\[2px] 0 
\end{pmatrix} \qquad
\begin{pmatrix}
  1 & \frac{1}{2} \\[2px]
  0 & \frac{1}{2} 
\end{pmatrix} 
\begin{pmatrix}
  0 \\[2px] 1
\end{pmatrix} = \begin{pmatrix}
\frac{1}{2} \\[2px] \frac{1}{2} 
\end{pmatrix} 
\] \vspace{5px}

All probabilistic operations are defined by \textit{stochastic} matrices, which
are matrices satisfying the following properties: 

\begin{mainbox}{Stochastic Matrices}
  \begin{itemize}
    \item[1.] All entries are nonnegative real numbers
    \item[2.] The entries in every column sum to 1
  \end{itemize}
\end{mainbox}

Equivalently, stochastic matrices are matrices whose columns all form
probability vectors

\subsection{Compositions of Probabilistic Operations}

Suppose $X$ is a system having classical state $\Sigma$, and $M_1, \cdots, M_n$
are stochastic matrices representing probabilistic operations on $X$. \\

If we apply $M_1$ to the probabilistic state represented by a probability
vector $u$, the resulting probabilistic state is $M_1 u.$ If we then apply
a second operation $M_2$, we obtain the probability vector

\[
M_2(M_1 u) = (M_2M_1)u. 
\] \vspace{5px}

The order in which operations are applied in a composition can change the
resulting operation. 

\section{Quantum Information}

Now we move onto quantum information where we use a different type of vector to
represent  \textit{quantum} states. In this section, we'll be considered with
systems that are finite in length with nonempty sets of classical states. We
will continue using \textit{Dirac} notation. 

\subsection{Quantum State Vectors} 

A \textit{quantum state} is represented the same way as probabilistic states.
Vectors representing quantum states are characterized by these following
properties: 

\begin{subbox}{Properties of Quantum States}
  \begin{itemize}
    \item[1.] The entries of a quantum state vector are \textit{complex \#s}
    \item[2.] The sum of the \textit{absolute values squared}  of the entries
      of a quantum state vector is 1. 
  \end{itemize}
\end{subbox}

Any speedup from a quantum computer, or improvement in communication protocol,
ultimately derives from these simple mathematical changes. \\

The \textit{Euclidean norm} of a column  vector 

\[
v = \begin{pmatrix}
  \alpha_1 \\ \vdots \\ \alpha_n
\end{pmatrix} 
\] \vspace{5px}

is defined as follows: 

\[
  || v || = \sqrt{\sum_{k=1}^n |\alpha_k|^2}
\] \vspace{5px}

Thus, quantum state vectors are  \textit{unit vectors} with respect to the
Euclidean norm. 

\paragraph{Examples of Qubit States} \mbox{} \\

The term \textit{qubit} refers to a quantum system whose classical state set is
{0, 1} -- the same as a bit, but a quantum state. 

\[
\begin{pmatrix}
  1 \\ 0 
\end{pmatrix} = |0\rangle \quad \text{and} \quad \begin{pmatrix}
  0 \\ 1
\end{pmatrix} = |1\rangle
\] \vspace{5px}
\[
\begin{pmatrix}
  \frac{1}{\sqrt{2}} \\[4px] \frac{1}{\sqrt{2}}
\end{pmatrix} = \frac{1}{\sqrt{2}}|0\rangle + \frac{1}{\sqrt{2}}|1\rangle
\]\vspace{5px}
and 

\[
\begin{pmatrix}
  \frac{1+2i}{3} \\[4px] -\frac{2}{3} 
\end{pmatrix} = \frac{1+2i}{3}|0\rangle - \frac{2}{3}|1\rangle
\] \vspace{5px}

Computing the sum of the absolute values squared of the entries respectively
yields

\[
|1|^2 + |0|^2 = 1 \quad \text{and}  \quad |0|^2 + |1|^2 = 1, 
\] \vspace{5px}

\[
  \Big| \frac{1}{\sqrt{2}} \Big|^2 + \Big| \frac{1}{\sqrt{2}}\Big|^2
  = \frac{1}{2} + \frac{1}{2} = 1
\] \vspace{5px}

and 

\[
\Big| \frac{1+2i}{3} \Big|^2 + \Big| -\frac{2}{3}\Big|^2 = \frac{5}{9}
+ \frac{4}{9} = 1
\] \vspace{5px}

These are therefore all valid quantum state vectors and are all linear
combinations, \textit{superpositions} of the 0 and 1 states. \\

The second example of a qubit state vector is commonly encountered and is
called the \textit{plus} state: 

\[
  |+\rangle = \frac{1}{\sqrt{2}}|0\rangle + \frac{1}{\sqrt{2}} |1\rangle
\] \vspace{5px}

We also use the notation 

\[
  |-\rangle = \frac{1}{\sqrt{2}}|0\rangle - \frac{1}{\sqrt{2}}|1\rangle
\] \vspace{5px}

for the \textit{minus} state. \\ 

It is common to use the notation $|\psi\rangle$ to denote an arbitrary vector
that may not necessarily be a standard basis vector. \\

If we have a vector $|\psi\rangle$ whose indices correspond to some classical
state set $\Sigma$, and if $a \in \Sigma$ is an element of this classical state
set, then the (matrix) product $\langle a | \psi \rangle $ is equal to the
entry of the vector  $|\psi\rangle$ whose index corresponds to $a$. \\ 

For example, if  $\Sigma = {0, 1}$ and 

\[
|\psi\rangle = \frac{1+2i}{3}|0\rangle - \frac{2}{3}|1\rangle = \begin{pmatrix}
  \frac{1+2i}{3}\\[4px] -\frac{2}{3}
\end{pmatrix} 
\] \vspace{5px}
then 

\[
\langle 0 | \psi \rangle = \frac{1+2i}{3} \quad \text{and}  \quad \langle
1 | \psi \rangle = -\frac{2}{3}
\] \vspace{5px}

Row vector notation, $\langle \psi |$ refers to the row vector obtained by
taking the  \textit{conjugate transpose} of the column vector $|\psi\rangle$.\\

The same $|\psi\rangle$ defined above would define $\langle \psi |$ 

\[
\langle \psi | = \frac{1-2i}{3}\langle 0| - \frac{2}{3}\langle 1| = \left(\frac{1-2i}{3} \quad -\frac{2}{3}\right)
\] \vspace{5px}

The reason we take the complex conjugate in addition to the transpose will
become clear when we discuss the \textit{inner product}.


\subsection{Measuring Quantum States}

The rule is simple: If a quantum state is measured, each classical state of the
system results with probability equal to the \textit{absolute value squared} of
the entry in the quantum state vector corresponding to that classical state.
This is known as the \textit{Born rule} in quantum mechanics. \\

For example, measuring the plus state 

\[
  |+\rangle = \frac{1}{\sqrt{2}}|0\rangle + \frac{1}{\sqrt{2}}|1\rangle
\] \vspace{5px}

results in two outcomes, 0 and 1, with probabilities 

\[
  \text{Pr(outcome is 0)} = |\langle 0 | +\rangle |^2
  = \Big|\frac{1}{\sqrt{2}}\Big|^2 = \frac{1}{2}
\] \vspace{5px}

\[
  \text{Pr(outcome is 1)} = |\langle 1 | + \rangle|^2
  = \Big|\frac{1}{\sqrt{2}}\Big|^2 = \frac{1}{2}
\] \vspace{5px}

\subsection{Unitary Operations}

The reason quantum information is fundamentally different from classical
information is because the set of allowable \textit{operations} that can be
performed on a quantum state is different than it is for classical information.
All operations on quantum states are defined with \textit{Unitary Matrices}.
A square matrix $U$ having complex \# entries is \textit{unitary} if it
satisfies the equations 

\begin{mainbox}{Unitary Matrices}
  \[
  UU^\dagger = I 
  \]
  \[
  U^\dagger U = I
  \]  
\end{mainbox}

Here, $I$ is the identity matrix, and $U^\dagger$ is the \textit{conjugate
transpose} of $U$. If $M$ is not a square matrix it is possible for $M^\dagger
M = I$ but not  $M M^\dagger = I$. This equivalence is only true for square
matrices. 

The condition that  $U$ is unitary is equivalent to the condition that
multiplication by $U$ does not change the euclidean norm of any vector. \\

In this sense, Unitary matrices are the quantum analogous of the classical
Stochastic matrices. 

\subsection{Important examples of Unitary Operations on Qubits}

\paragraph{\textit{Pauli Operations}} \mbox{} \\

The four Pauli matrices are as follows: 

\[
I = \begin{pmatrix}
  1 & 0 \\ 0 & 1
\end{pmatrix}, \quad \sigma_x = \begin{pmatrix}
  0 & 1 \\[2px]
  1 & 0 
\end{pmatrix}, \quad \sigma_y = \begin{pmatrix}
  0 & -i \\[2px]
  i & 0 
\end{pmatrix}, \quad \sigma_z = \begin{pmatrix}
  1 & 0 \\[2px]
  0 & -1 
\end{pmatrix} 
\] \vspace{5px}

A common notation is that $X = \sigma_x$, $Y = \sigma_y$, and $Z = \sigma_z$.
The  $X$ operation is also called the \textit{bit flip} or a  \textit{NOT}
operation because it induces this action on bits: 

 \[
X|0\rangle = |1\rangle \quad \text{and} \quad X|1\rangle = |0\rangle
\] \vspace{5px}

The $Z$ operation is also called a \textit{phase flip} because it has this
action: 

\[
Z|0\rangle = |0\rangle \quad \text{and} \quad Z|1\rangle = -|1\rangle
\] \vspace{5px}

\paragraph{\textit{Hadamard Operation}} \mbox{} \\ 

The Hadamard Operation is described by this matrix: 

\[
H = \begin{pmatrix}
  \frac{1}{\sqrt{2}} & \frac{1}{\sqrt{2}} \\[4px]
  \frac{1}{\sqrt{2}} & -\frac{1}{\sqrt{2}} 
\end{pmatrix}
\] \vspace{5px}

Applying the Hadamard Gates on well-known quantum states yields 

\begin{align*}
  H|0\rangle &= |+\rangle \\
  H|1\rangle &= |-\rangle \\
  H|+\rangle &= |0\rangle \\
  H|-\rangle &= |1\rangle
\end{align*}

It is worth pausing to consider the fact that $H|+\rangle = |0\rangle$ and
$H|-\rangle = |1\rangle$. Both  $|+\rangle$ and  $|-\rangle$ produce
0 and 1 with the same probabilities as each other -- so one would think having
both states is useless and we wouldn't know which of the states a qubit is in 
-- $|+\rangle$ or $|-\rangle$. However, if we apply a Hadamard Operation and
then measure we obtain 0 with certainty if the original state was $|+\rangle$ 
and 1 if the original state was $|-\rangle$. \\ 

  Thus the quantum states $|+\rangle$ and $|-\rangle$ can be discriminated
  \textit{perfectly}. This reveals that sign changes, or more generally changes
  to the \textit{phases} of the complex \# entries of a quantum state vector,
  an significantly change the state.\\

  Here is another example, this time of the action of a $T$ operation on a plus
  state: 

  \[
    T|+\rangle = T\left( \frac{1}{\sqrt{2}} |0\rangle + \frac{1}{\sqrt{2}}|1\rangle \right) = \frac{1}{\sqrt{2}}T|0\rangle + \frac{1}{\sqrt{2}}|1\rangle = \frac{1}{\sqrt{2}} |0\rangle + \frac{1+i}{2}|1\rangle
  \] \vspace{5px}
  
Notice we did not bother to convert to the equivalent matrix/vector forms. We instead used the linearity of matrix multiplication together with the formulas

\[
  T|0\rangle = |0\rangle \quad \text{and} \quad T|1\rangle = \frac{1+i}{\sqrt{2}}|1\rangle
\] \vspace{5px}



\paragraph{\textit{Phase Operations}} \mbox{} \\

A phase operation is one described by the matrix 

\[
P_\theta = \begin{pmatrix}
  1 & 0 \\[2px]
  0 & e^{i\theta} 
\end{pmatrix} 
\] \vspace{5px}

for any choice of real number $\theta$, that is, any complex \# with a norm of
1. The operations 

\[
  S = P_{\pi / 2} = \begin{pmatrix}
    1 & 0 \\[4px]
    0 & i 
  \end{pmatrix} \quad \text{and}\quad T = P_{\pi / 4} = \begin{pmatrix}
    1 & 0 \\[4px]
    0 & \frac{1+i}{\sqrt{2}} 
  \end{pmatrix}  \] \vspace{5px}

are particularly important examples. Other examples include $I = P_0$ and $Z
= P_\pi$

\subsection{Compositions of Qubit Unitary Operations} 

If we first apply a Hadamard Operation, followed by the $S$ operation, followed by another Hadamard Operation, the resulting operation $R$ is as follows: 

\[
R = HSH = \begin{pmatrix}
  \frac{1}{\sqrt{2}} & \frac{1}{\sqrt{2}} \\[4px]
  \frac{1}{\sqrt{2}} & -\frac{1}{\sqrt{2}} 
\end{pmatrix} \begin{pmatrix}
  1 & 0 \\[4px]
  0 & i 
\end{pmatrix} \begin{pmatrix}
  \frac{1}{\sqrt{2}} & \frac{1}{\sqrt{2}} \\[4px]
  \frac{1}{\sqrt{2}} & -\frac{1}{\sqrt{2}} 
\end{pmatrix}  = \begin{pmatrix}
  \frac{1+i}{2} & \frac{1-i}{2} \\[4px]
  \frac{1-i}{2} & \frac{1+i}{2} 
\end{pmatrix} 
\] \vspace{5px}

This unitary operation $R$ is an interesting example. By applying this operation twice, we obtain a \textit{NOT} operation. 

\[
R^2 = \begin{pmatrix}
  \frac{1+i}{2} & \frac{1-i}{2} \\[4px]
  \frac{1-i}{2} & \frac{1+i}{2} 
\end{pmatrix}^2 = \begin{pmatrix}
  0 & 1 \\[4px]
  1 & 0 
\end{pmatrix} 
\] 
\section{Qiskit Examples}

In python, matrix and vector computations can be performed using the \verb|array| class from the \verb|NumPy| library. Here is an example of a code cell that defines two vectors, \verb|ket0| and \verb|ket1| corresponding to their qubit state vectors $|0\rangle$ and $|1\rangle$, and displays their average.  \\

\begin{lstlisting}
from numpy import array

ket_0 = array([1, 0])
ket_1 = array([0, 1])

ket_0 / 2 + ket_1 / 2 
\end{lstlisting} \vspace{5px}

Output:  

\begin{lstlisting}
array([0.5, 0.5])
\end{lstlisting} \vspace{5px}


We can also use \verb|array| to create matrices that represent operations.\\  

\begin{lstlisting}
M_1 = array([[1, 1], [0, 0]])
M_2 = array([[1, 1], [1, 0]])

M_1 / 2 + M_2 / 2
\end{lstlisting}

Output: 

\begin{lstlisting}
array([[1. , 1. ],
       [0.5, 0. ]])
\end{lstlisting} \vspace{5px}

Matrix multiplication can be performed using \verb|matmul| function from \verb|NumPy| : \\

\begin{lstlisting}
from numpy import matmul 

display(matmul(M_1, ket_1))
display(matmul(M_1, M_2))
display(matmul(M_2, M_1))
\end{lstlisting} \vspace{5px}

Output: 

\begin{lstlisting}
array([1, 0])
array([[2, 1], 
       [0, 0]])
array([[1, 1], 
       [1, 1]])
\end{lstlisting} \vspace{5px}



































\end{document}

